\section{From Lesson 2 to Lab 2}
\subsection{Bounded help}

\begin{frame}{Give the agent one task you can evaluate}
  A broad request hides the important decisions:
  \begin{block}{Too broad}
    Clean this dataset and make it analysis-ready.
  \end{block}

  Ask for one improvement:
  \begin{itemize}
    \item add a key or range assertion;
    \item diagnose unmatched join keys;
    \item produce a compact build record.
  \end{itemize}

  \vfill
  First write a small pipeline that you can read. Then use the agent as a reviewer.
\end{frame}

\begin{frame}[fragile]{A bounded request for the health table}
\begin{lstlisting}
Goal: Add a build record for the health analysis table.
Context: Read code/02-build-analysis.R and the indicator dictionary.
Permission: Propose a diff. Do not run it or change frozen data.
Constraints: Keep the country-year unit and lowercase names.
Proof: Report rows, unique keys, years, missingness, ranges,
unmatched joins, and every failed condition.
\end{lstlisting}
\end{frame}

\begin{frame}{Review the proposed diff}
  For every added check, answer:
  \begin{enumerate}
    \item Which table-contract field does it test?
    \item What failure would it expose?
    \item Can I explain the output?
    \item Does the proposed change affect the population, transformation, or key?
  \end{enumerate}

  \vfill
  A change to the population, transformation, or key is a substantive decision.
  The analyst remains responsible even when the request is described as code cleanup.
\end{frame}

\begin{frame}{Decisions the analyst keeps}
  \begin{itemize}
    \item The meaning of one row.
    \item The population represented by a filter.
    \item Whether a transformation fits the policy question.
    \item Whether current metadata can be interpreted historically.
    \item Whether missing coverage weakens the comparison.
  \end{itemize}

  \vfill
  The analyst interprets the evidence that the agent exposes.
\end{frame}

\begin{frame}{Updating public data is an extension}
  Extending one indicator beyond 2022 can be useful. It also adds:
  \begin{itemize}
    \item an internet and API dependency;
    \item retrieval dates and version changes;
    \item possible revisions to earlier years;
    \item new metadata questions.
  \end{itemize}

  \vfill
  First build the frozen table correctly. Then extend it with a separate record.
\end{frame}

\begin{frame}{What we learned in Lesson 2}
  \begin{itemize}
    \item The table contract comes before the data verbs.
    \item Every verb should have a reason and a check.
    \item A join must preserve the intended key and unit.
    \item The analysis table should rebuild from a protected source.
    \item The agent improves one bounded part under analyst review.
  \end{itemize}
\end{frame}

\begin{frame}{Lab 2}
  \vfill
  \begin{center}
    {\LARGE\color{crimson}Build one country--year table.}

    \vspace{0.75em}
    Test its contract and explain the decisions that produced it.
  \end{center}
  \vfill
\end{frame}

\begin{frame}{Lab 2 route}
  \small
  \begin{tabularx}{\textwidth}{@{}>{\color{crimson}\bfseries}p{1.6cm}>{\bfseries}p{3.2cm}X@{}}
    \toprule
    Minutes & Work & Result \\
    \midrule
    00--10 & Contract & Population, unit, key, period, and variables. \\
    10--25 & Select + filter & Required columns and sample. \\
    25--35 & Mutate + group & One justified variable. \\
    35--45 & Join & Metadata after a key check. \\
    45--52 & Agent pass & One reviewable improvement. \\
    52--60 & Build record & Checks and a short explanation. \\
    \bottomrule
  \end{tabularx}
\end{frame}

\begin{frame}{The lab is done when}
  \begin{itemize}
    \item The script rebuilds the table from the frozen file.
    \item The country--year key is unique.
    \item The build record includes rows, years, missingness, ranges, and matches.
    \item You can explain one consequential decision without the agent.
  \end{itemize}

  \vfill
  The lab has no submission or grade. Keep the script because the next lesson
  will use this table.
\end{frame}
